\begin{ExerciseBlock}
\ExerciseHeading{02}{Gedämpfte Schwingung}{free_oscillations_damping_from_amplitudes}
Die Maximalausschläge eines Federpendels nehmen mit der Zeit ab. Zu Beginn ($t=0$) beträgt die Amplitude $x_0=\qty{5.0}{\centi\metre}$. In den nächsten \qty{10}{\second} führt das Pendel genau fünf Schwingungen aus. Danach beträgt die Amplitude $x_5=\qty{0.5}{\centi\metre}$.
\begin{enumerate}
\item Berechnen Sie $T$, $\omega$, die Abklingkonstante $\delta$ und die ungedämpfte Kreisfrequenz $\omega_0$.
\item Um wie viel Prozent ist $\omega_0$ größer als $\omega$?
\end{enumerate}
\end{ExerciseBlock}
\UseExerciseSolution{free_oscillations_damping_from_amplitudes}
